
\begin{frame}
\frametitle{Plan du cours}

\alert{Tout} ce qu'il faut comprendre pour \alert{utiliser} efficacement un SGBD relationnel.

\begin{redbullet}
  \item Architecture 
  \item Structuration d'une base relationnelle
  \item Interrogation: approches déclarative et procédurale
  \item Conception
  \item Quelques aspects pratiques (schémas, programmation)
  \item Transactions
\end{redbullet}

\vfill

\begin{blocgauche}
Pas (ou très peu) d'information sur le fonctionnement interne du système
\end{blocgauche}

\end{frame}



\begin{frame}{Organisation du cours}

Le cours est découpé en \alert{chapitres}, couvrant un sujet bien déterminé, et en \alert{sessions}.

\smallskip

\alert{L'unité de travail est la session}. Chaque session
demande environ 2 ou 3  heures de travail personnel (bien sûr, cela dépend également de vous).

\vfill


Pour assimiler une session vous pouvez combiner les ressources suivantes:

\begin{redbullet}
  \item \alert{La lecture} du \href{http://sql.bdpedia.fr}{support en ligne}, 
       également disponible en \href{http://sql.bdpedia.fr/files/cbd-sql.pdf}{PDF} ou en 
       \href{http://sql.bdpedia.fr/files/cbd-sql.epub}{ePub}.
    \item \alert{Le suivi du cours}, en vidéo ou en présentiel.
    \item Vous devriez alors  savoir répondre aux Quiz en fin de session
   \item \alert{La réalisation des exercices et des ateliers} proposés en fin de chapitre.
\end{redbullet}

\vfill
\begin{blocgauche}
\alert{La réalisation des exercices est essentielle} pour vérifier que vous maîtrisez le
contenu.
\end{blocgauche}

\end{frame}


\begin{frame}{Les données}

\alert{Donnée} = valeur numérisée décrivant de manière élémentaire
un fait, une mesure, une réalité

\smallskip

\alert{Exemple}: le nom de l'auteur, l'âge du capitaine, le titre du livre ...


\vfill

Les données décrivent des \alert{entités} du monde réel, elles-mêmes associées les unes aux autres.

\smallskip
\alert{Exemple}: \textit{Nicolas Bouvier est un écrivain suisse auteur de récit de voyage
culte "l'usage du monde" paru en 1963}: deux entités, liées par la notion d'auteur.

\vfill
\begin{blocgauche}
Une base de données a donc une \alert{structure}, sinon c'est autre chose
(une collection, un tas de documents, textes ou images).
\end{blocgauche}

\end{frame}


\begin{frame}{Base de  données}

Caractères essentiels: structuration, persistance (vs. volatilité).

\vfill

\begin{defblock}{Base de données}
Une base de données est un ensemble d’informations structurées mémorisées sur un support persistant.
\end{defblock}

\vfill


\begin{blocgauche}
Des fichiers structurés (tableur, CSV) sont des bases de données.
\end{blocgauche}

\end{frame}


\begin{frame}[fragile]{Exemple de fichiers structurés}

Format CSV: une ligne par entité ; champs séparés par des ';'
\begin{small}
\begin{minted}{C}
"Bouvier" ; "Nicolas"; "L'usage du monde" ; 1963
\end{minted}
\end{small}
\vfill

Base de données = 2, 10 ou 1 million de lignes sur le même format.

\begin{small}
\begin{minted}{C}
"Bouvier"   ; "Nicolas"; "L'usage du monde" ; 1963
"Stevenson" ; "Robert-Louis"  ; "Voyage dans les Cévennes avec un âne" ; 1879
...
\end{minted}
\end{small}

\vfill
Suffisant?
\end{frame}

%%%%%%%%%%%%%%%
\begin{frame}{Fichiers = base de données ? }

Peut-on construire des applications directement sur des fichiers?
\vfill

\figSlideWithSize{sans-serveur}{7cm}

\vfill

Insurmontables problèmes de productivité, de fiabilité, d'efficacité
\end{frame}



%%%%%%%%%%%%%%%
\begin{frame}{Les SGBD }

Système informatique qui assure la gestion de l'ensemble des informations
stockées dans une base de données.


\vfill

\figSlideWithSize{avec-serveur}{10cm}

\end{frame}


%%%%%%%%%%%%%%%
\begin{frame}{Niveaux d'abstraction et modèle de données}

Le serveur peut présenter une représentation \alert{logique} des données très éloignée
de la représentation \alert{physique}.

\vfill

\figSlideWithSize{physique-logique}{10cm}

\vfill

Le niveau logique définit la \alert{modélisation} des données.


\end{frame}


%%%%%%%%%%%%%%%
\begin{frame}{Langages - SQL}

Pour manipuler les données, le serveur propose un \alert{langage d'interrogation}.

\vfill

Le plus répandu est SQL: interrogation, mise à jour, contraintes sur la base, droits d'accès...

\vfill

\figSlideWithSize{physique-logique-sql}{9.5cm}

\end{frame}


\begin{frame}{À retenir}

Les \alert{bases de données} sont des ensembles structurés stockés dans des fichiers.

\vfill 

Les \alert{systèmes de gestion de base de données (SGBD)} sont des systèmes qui prennent en charge
toute la complexité de gestion des fichiers.

\vfill

\vfill 

Un SGBD propose aux applications une vue \alert{logique} indépendante du niveau \alert{physique} (stockage).

\vfill
\begin{blocgauche}
SQL est le langage qui permet d'interagir avec la vue logique d'une base de données (relationnelle)
\end{blocgauche}
\end{frame}

